% GENERATED FILE. Edit canonical lessons, then run npm run generate:books.
% canonical-source-hash: fnv1a64:e5007f8779db0368
% canonical-lessons: LA-C350-obsono, LA-C350-obsonium, LA-C350-abluo, LA-C350-tergeo, LA-C350-cultellus

\chapter{In macellō}
\label{ch:la-in-macello}
\hlchaptermodality{350}

\begin{chapteropening}
\textbf{By the end of this chapter:} \emph{I can tell how food was bought and made ready: to buy food, the provisions, to wash off, to wipe clean and the small knife.}

Complete the last lesson of chapter 350: I can tell how food was bought and made ready: to buy food, the provisions, to wash off, to wipe clean and the small knife.
\end{chapteropening}

\section[obsōnō, obsōnāre]{obsōnō, obsōnāre — to buy food, to shop for a meal}

\label{lesson:LA-C350-obsono}



\begin{quote}
Before the new one: say the Latin for unbelievable, then the Latin for an invention.
\end{quote}

\subsection*{You'll want to know: obsōnō, obsōnāre}
\textbf{obsōnō, obsōnāre} — \textquotedblleft{}to buy food, to shop for a meal\textquotedblright{}.

\begin{quote}
Coquus māne in macellō obsōnābat.
\end{quote}

\begin{quote}
\emph{In the morning the cook was buying food at the market.}
\end{quote}

\emph{obsōnō, obsōnāre, obsōnāvī, obsōnātum}, also \emph{opsōnō}. Comedies are full of cooks sent off to shop.

\begin{cousinweb}
\textbf{obsōnō} was borrowed from Greek \emph{opsōneō}, \textquotedblleft{}to buy fish or other dainties\textquotedblright{} (EtymDB, Lewis \& Short).
\end{cousinweb}

\begin{grammarlens}[title={What you've built}]
One more word for telling, explaining and describing at B1.
\end{grammarlens}

\subsection*{Your turn}
\begin{itemize}
\raggedright
  \item \emph{Say it:} \emph{obsōnō, obsōnāre}
  \item \emph{Say it:} \emph{obsōnō, obsōnāre}, once more
  \item \emph{Say it:} a sentence of your own with \emph{obsōnō, obsōnāre}
  \item \emph{Recall:} say the Latin for unbelievable, then the Latin for an invention, then say \emph{obsōnō, obsōnāre} again
\end{itemize}

\begin{tcolorbox}[breakable,colback=teal!4,colframe=teal!35!black,title={Before you move on}]
What does obsōnō, obsōnāre mean? (\textquotedblleft{}To buy food, to shop for a meal\textquotedblright{}.) Say it once more.
\end{tcolorbox}
\section[obsōnium, obsōniī]{obsōnium, obsōniī — food to go with bread, provisions}

\label{lesson:LA-C350-obsonium}



\begin{quote}
Before the new one: say the Latin for an invention, then the Latin for to buy food.
\end{quote}

\subsection*{You'll want to know: obsōnium, obsōniī}
\textbf{obsōnium, obsōniī} — \textquotedblleft{}food to go with bread, provisions\textquotedblright{}.

\begin{quote}
Obsōnium in culīnam portāvit: piscēs et holera.
\end{quote}

\begin{quote}
\emph{He carried the provisions into the kitchen: fish and vegetables.}
\end{quote}

Neuter. Anything eaten with bread, very often fish.

\begin{cousinweb}
\textbf{obsōnium} was borrowed from Greek \emph{opsōnion}, \textquotedblleft{}provisions, what is eaten with bread\textquotedblright{} (EtymDB).
\end{cousinweb}

\begin{grammarlens}[title={What you've built}]
One more word for telling, explaining and describing at B1.
\end{grammarlens}

\subsection*{Your turn}
\begin{itemize}
\raggedright
  \item \emph{Say it:} \emph{obsōnium, obsōniī}
  \item \emph{Say it:} \emph{obsōnium, obsōniī}, once more
  \item \emph{Say it:} a sentence of your own with \emph{obsōnium, obsōniī}
  \item \emph{Recall:} say the Latin for an invention, then the Latin for to buy food, then say \emph{obsōnium, obsōniī} again
\end{itemize}

\begin{tcolorbox}[breakable,colback=teal!4,colframe=teal!35!black,title={Before you move on}]
What does obsōnium, obsōniī mean? (\textquotedblleft{}Food to go with bread, provisions\textquotedblright{}.) Say it once more.
\end{tcolorbox}
\section[abluō, abluere]{abluō, abluere — to wash off, to wash clean}

\label{lesson:LA-C350-abluo}



\begin{quote}
Before the new one: say the Latin for to buy food, then the Latin for food to go with bread.
\end{quote}

\subsection*{You'll want to know: abluō, abluere}
\textbf{abluō, abluere} — \textquotedblleft{}to wash off, to wash clean\textquotedblright{}.

\begin{quote}
Coquus holera aquā frīgidā abluit.
\end{quote}

\begin{quote}
\emph{The cook washed the vegetables in cold water.}
\end{quote}

\emph{abluō, abluere, abluī, ablūtum}. English \emph{ablution}.

\begin{cousinweb}
\textbf{abluō} joins \emph{ab}, \textquotedblleft{}away\textquotedblright{}, and \emph{luō}, \textquotedblleft{}to wash\textquotedblright{}, a cousin of \emph{lavō} (EtymDB). English \emph{ablution} comes from it.
\end{cousinweb}

\begin{grammarlens}[title={What you've built}]
One more word for telling, explaining and describing at B1.
\end{grammarlens}

\subsection*{Your turn}
\begin{itemize}
\raggedright
  \item \emph{Say it:} \emph{abluō, abluere}
  \item \emph{Say it:} \emph{abluō, abluere}, once more
  \item \emph{Say it:} a sentence of your own with \emph{abluō, abluere}
  \item \emph{Recall:} say the Latin for to buy food, then the Latin for food to go with bread, then say \emph{abluō, abluere} again
\end{itemize}

\begin{tcolorbox}[breakable,colback=teal!4,colframe=teal!35!black,title={Before you move on}]
What does abluō, abluere mean? (\textquotedblleft{}To wash off, to wash clean\textquotedblright{}.) Say it once more.
\end{tcolorbox}
\section[tergeō, tergēre]{tergeō, tergēre — to wipe, to wipe clean}

\label{lesson:LA-C350-tergeo}



\begin{quote}
Before the new one: say the Latin for food to go with bread, then the Latin for to wash off.
\end{quote}

\subsection*{You'll want to know: tergeō, tergēre}
\textbf{tergeō, tergēre} — \textquotedblleft{}to wipe, to wipe clean\textquotedblright{}.

\begin{quote}
Ancilla mēnsam linteō tersit.
\end{quote}

\begin{quote}
\emph{The maid wiped the dining table with a cloth.}
\end{quote}

\emph{tergeō, tergēre, tersī, tersum}. Its compound \emph{dētergeō} gives English \emph{detergent}.

\begin{cousinweb}
\textbf{tergeō} is an inherited verb (EtymDB). English \emph{detergent} comes from its compound \emph{dētergeō}, \textquotedblleft{}to wipe off\textquotedblright{}.
\end{cousinweb}

\begin{grammarlens}[title={What you've built}]
One more word for telling, explaining and describing at B1.
\end{grammarlens}

\subsection*{Your turn}
\begin{itemize}
\raggedright
  \item \emph{Say it:} \emph{tergeō, tergēre}
  \item \emph{Say it:} \emph{tergeō, tergēre}, once more
  \item \emph{Say it:} a sentence of your own with \emph{tergeō, tergēre}
  \item \emph{Recall:} say the Latin for food to go with bread, then the Latin for to wash off, then say \emph{tergeō, tergēre} again
\end{itemize}

\begin{tcolorbox}[breakable,colback=teal!4,colframe=teal!35!black,title={Before you move on}]
What does tergeō, tergēre mean? (\textquotedblleft{}To wipe, to wipe clean\textquotedblright{}.) Say it once more.
\end{tcolorbox}
\section[cultellus, cultellī]{cultellus, cultellī — a small knife}

\label{lesson:LA-C350-cultellus}



\begin{quote}
Before the new one: say the Latin for to wash off, then the Latin for to wipe.
\end{quote}

\subsection*{You'll want to know: cultellus, cultellī}
\textbf{cultellus, cultellī} — \textquotedblleft{}a small knife\textquotedblright{}.

\begin{quote}
Coquus cultellum acūtum ē cistā sūmpsit.
\end{quote}

\begin{quote}
\emph{The cook took a sharp little knife from the box.}
\end{quote}

Masculine, a diminutive of \emph{culter} (taught). English \emph{cutlass} and \emph{cutlery} come from this family.

\begin{cousinweb}
\textbf{cultellus} is a diminutive of \emph{culter}, \textquotedblleft{}knife\textquotedblright{} (Lewis \& Short). English \emph{cutlass} and \emph{cutlery} come from it through French.
\end{cousinweb}

\begin{grammarlens}[title={What you've built}]
One more word for telling, explaining and describing at B1.
\end{grammarlens}

\subsection*{Your turn}
\begin{itemize}
\raggedright
  \item \emph{Say it:} \emph{cultellus, cultellī}
  \item \emph{Say it:} \emph{cultellus, cultellī}, once more
  \item \emph{Say it:} a sentence of your own with \emph{cultellus, cultellī}
  \item \emph{Recall:} say the Latin for to wash off, then the Latin for to wipe, then say \emph{cultellus, cultellī} again
\end{itemize}

\begin{tcolorbox}[breakable,colback=teal!4,colframe=teal!35!black,title={Before you move on}]
What does cultellus, cultellī mean? (\textquotedblleft{}A small knife\textquotedblright{}.) Say it once more.
\end{tcolorbox}
